\section{Curves}

\begin{theorem}[Euler-characteristic components]
\label{akcp-8-1-curve-picard}
\dcref{8.1}
\source{MR0555258-compactifying-picard:page-0052}
\uses{akcp-3-5-curve-sheaves,akcp-6-3-picard-integral}
If $X/S$ is flat locally projective finitely presented with geometrically
integral curve fibers, then
\[
\Pic^{\mathrm{tf},\acute{\mathrm et}}_{X/S}=\coprod_{n\in\mathbb Z}P_n,
\]
where $P_n$ parametrizes rank-one torsion-free sheaves of Euler characteristic
$n$ on the fibers.
\end{theorem}
\begin{proof}
Work locally with $X/S$ projective and connected.  The representability theorem
decomposes the Picard sheaf by Hilbert polynomial.  The degree of the fibers is
constant on a connected base, and \cref{akcp-3-5-curve-sheaves} gives
$\chi(\II(m))=\deg(X_s)m+\chi(\II)$, so polynomial components are exactly the
Euler-characteristic components.  Glue over the base.
\end{proof}

\begin{definition}[Components of the Abel map]
\label{akcp-8-2-abel-components}
\dcref{8.2}
\source{MR0555258-compactifying-picard:page-0052}
\uses{akcp-5-15-quot-strata-representable,akcp-5-16-abel-map,akcp-8-1-curve-picard}
For a flat locally projective family of integral curves and a relatively
torsion-free rank-one $\FF$, every nonzero pseudo-ideal is again torsion-free
rank one and Cohen--Macaulay.  Hence, for $d>0,
P\operatorname{-div}^{d}(\FF)=\Quot^{d}_{\FF/X/S}$, and the Abel map splits
into components
\[
\alpha_d:\Quot^{d}_{\FF/X/S}\longrightarrow P_{\chi(\FF_s)-d}.
\]
Tensoring by any invertible $\LL$ gives a commutative diagram identifying the
components for $\FF$ and $\FF\otimes\LL$.  For a Gorenstein family of curves,
tensoring by $\omega_{X/S}$ identifies the component for $\OO_X$ with that for
$\omega$.
\end{definition}

\begin{definition}[Index of specialty]
\label{akcp-8-3-index-specialty}
\dcref{8.3}
\source{MR0555258-compactifying-picard:page-0053}
\uses{akcp-3-5-curve-sheaves,akcp-8-2-abel-components}
For a projective integral curve, the $\FF$-index of specialty of a rank-one
torsion-free $\II$ is $\dim\Ext^1_X(\II,\FF)$.  For a nontrivial quotient
$\GG$ of $\FF$, it is $\dim\Ext^1_X(\II(\GG),\FF)$.  Index zero means
non-special.  Tensoring by an invertible sheaf preserves the index, and the
definition extends fiberwise and to scheme points of the Quot and Picard
spaces.  For $\FF=\omega$, duality identifies the index of a quotient with
$h^0(\II(\GG))$.
\end{definition}

\begin{theorem}[The Abel map for integral curves]
\label{akcp-8-4-abel-curve}
\dcref{8.4}
\source{MR0555258-compactifying-picard:page-0054}
\source{MR0555258-compactifying-picard:page-0055}
\uses{akcp-3-5-curve-sheaves,akcp-5-18-abel-smooth,akcp-5-20-abel-proper,akcp-8-1-curve-picard,akcp-8-2-abel-components,akcp-8-3-index-specialty}
Let $X/S$ be flat, finitely presented and locally projective with integral
curve fibers of common arithmetic genus $p$, and let $\omega$ be dualizing.
For the $d$th Abel component
\[
\alpha_d:\Quot^d_{\omega/X/S}\longrightarrow
P_{p-1-d}=\Pic^{\mathrm{tf}}_{X/S}(\chi=p-1-d),
\]
the following hold:
\begin{enumerate}
\item $\alpha_d$ is surjective iff $d\ge p$; if $d<p$ its image omits a point.
\item Over an $\omega$-nonspecial point it is smooth of relative dimension
      $d-p$.
\item If a neighborhood of a point in the closure of the invertible Picard
      locus has nonempty fibers of constant dimension, then $d\ge p$ and all
      points there are $\omega$-nonspecial.
\item A flat point of $\alpha_d$ lying over that closure is nonspecial.
\item $\alpha_d$ is smooth iff $d\ge2p-1$.
\end{enumerate}
\end{theorem}
\begin{proof}
For a geometric point represented by $\II$, the fiber dimension is
\[
\dim\Hom(\II,\omega_s)-1=h^1(\II)-1
\]
by \cref{akcp-5-18-abel-smooth} and duality.  If $d\ge p$, then
$\chi(\II)\le0$, so $h^1(\II)>0$ and the Abel fiber is nonempty.  If $d<p$,
choose an invertible line bundle with $h^1=0$ from
\cref{akcp-3-5-curve-sheaves}; its Picard point is omitted.  Nonspecial fibers
have the stated smoothness and dimension by \cref{akcp-5-18-abel-smooth}.

For (3), over an algebraically closed field the invertible Picard component is
irreducible and the proper image of $\alpha_d$ is closed.  If it contains a
nonempty open subset it contains the whole component, forcing $d\ge p$ by (1).
The nonspecial locus is open by upper semicontinuity.  It meets every relevant
fiber by the line-bundle constructions in \cref{akcp-3-5-curve-sheaves}; hence
irreducibility makes it meet the constant-dimension neighborhood, proving (3).
For (4), flatness makes the image of a connected neighborhood open with
constant-dimensional fibers; apply (3).  If $d\ge2p-1$, then
$\chi(\II)<-p$ and \cref{akcp-3-5-curve-sheaves} gives $h^0(\II)=0$, so every
fiber is nonspecial and (2) gives smoothness.  Conversely, for
$d\le2p-2$, the sheaves constructed in (3.5) have different $h^1$, hence the
fiber dimensions differ; the Abel map cannot be smooth.
\end{proof}

\begin{theorem}[Projectivity of fixed Euler components]
\label{akcp-8-5-fixed-euler}
\dcref{8.5}
\source{MR0555258-compactifying-picard:page-0056}
\uses{akcp-6-6-picard-cm,akcp-8-1-curve-picard,akcp-8-4-abel-curve}
For a flat locally projective family of integral curves of genus $p$, each
$P_n$ is finitely presented and locally quasi-projective; if $X/S$ is
projective, it is projective; if $S$ is connected, it is strongly projective.
\end{theorem}
\begin{proof}
The Hilbert polynomial of a rank-one sheaf is
$\deg(X_s)m+n$, with constant leading coefficient on a connected base, so
$P_n$ is strongly quasi-projective by the scheme representation theorem.  Choose
$m$ with the corresponding polynomial negative and twist by $\OO(m)$, reducing
to a sufficiently large Abel component.  The Abel map from the projective Quot
scheme is surjective by \cref{akcp-8-4-abel-curve}, so the component is
universally closed.  Its inclusion into an ambient projective bundle is then
closed, hence projective and strongly projective.
\end{proof}

\begin{theorem}[D'Souza--Rego criterion]
\label{akcp-8-6-dsouza-rego}
\dcref{8.6}
\source{MR0555258-compactifying-picard:page-0057}
\source{MR0555258-compactifying-picard:page-0058}
\uses{akcp-1-10-open-ext,akcp-3-5-curve-sheaves,akcp-5-18-abel-smooth,akcp-8-4-abel-curve}
For a flat finitely presented locally projective family of integral curves of
genus $p$ and $d>0$, the following are equivalent for
$\alpha_d:\Hilb^d_{X/S}\to P_{1-p-d}$:
\begin{enumerate}
\item $d\ge2p-1$ and every fiber is Gorenstein;
\item $\alpha_d$ is smooth of relative dimension $d-p$;
\item all fibers of $\alpha_d$ have the same dimension;
\item $d\ge2p-1$ and fibers over the closure of the invertible Picard locus
      have the same dimension.
\end{enumerate}
\end{theorem}
\begin{proof}
The implication (1)$\Rightarrow$(2) is \cref{akcp-8-4-abel-curve}, and
(2)$\Rightarrow$(3) is immediate.  Over an algebraically closed field, the
line-bundle examples in \cref{akcp-3-5-curve-sheaves} give differing $h^1$
for $d\le2p-2$; \cref{akcp-5-18-abel-smooth} then gives differing fiber
dimensions, proving (3)$\Rightarrow$(4).

Assume (4).  Put $L$ of degree $d-1$ and examine the diagonal exact sequence
on $X\times X$:
\[
0\to\OO_{\Delta}\to\OO_{X\times X}\to\II_\Delta\to0.
\]
Applying $\mathcal H\!om(-,L\boxtimes M)$ and pushing forward along the first
projection yields an exact sequence relating $f_*L\otimes M$, the Hom sheaf
of $\II_\Delta$, and $\Ext^1(\II_\Delta,L\boxtimes M)$.  If $h^0(L)\ne0$,
the last assertion of \cref{akcp-3-5-curve-sheaves} identifies $L$ with the
dualizing sheaf, so every fiber is Gorenstein.  Otherwise exchange
(
\cref{akcp-1-10-open-ext}) makes $M\mapsto f_*(L\otimes M)$ exact.  Constant
fiber dimensions and \cref{akcp-5-18-abel-smooth} make
$\mathcal H\!om(\II_\Delta,L\boxtimes\OO_X)$ locally free; the preceding exact
sequence then makes
$M\mapsto f_*\Ext^1(\II_\Delta,L\boxtimes M)$ exact.  Since this Ext sheaf is
supported on the diagonal, it is flat over $X$ and commutes with base change.
Its fiber at a closed point is $\Ext^1(k(x),\OO_{X})$, a skyscraper of rank one
at smooth points; constant rank forces rank one everywhere.  Thus every fiber is
Gorenstein, and (1) follows together with $d\ge2p-1$ from (4).
\end{proof}

\begin{lemma}[Hilbert scheme of the diagonal]
\label{akcp-8-7-diagonal-hilbert}
\dcref{8.7}
\source{MR0555258-compactifying-picard:page-0059}
The diagonal gives an isomorphism $X\xrightarrow{\sim}\Hilb^1_{X/S}$ for a
flat finitely presented locally quasi-projective family.
\end{lemma}
\begin{proof}
The diagonal is a flat family with fiber Hilbert polynomial $1$.  Conversely,
a member $Y\subset X_T$ has every fiber equal to one point, hence
$Y\to T$ is a surjective closed immersion.  A flat finitely presented
surjective closed immersion is an isomorphism (its ideal vanishes fiberwise and
then by Nakayama), so $Y$ is the graph of a unique $T\to X$.
\end{proof}

\begin{theorem}[The first Abel component]
\label{akcp-8-8-first-abel}
\dcref{8.8}
\source{MR0555258-compactifying-picard:page-0059}
\source{MR0555258-compactifying-picard:page-0060}
\uses{akcp-5-17-abel-fiber,akcp-5-20-abel-proper,akcp-8-4-abel-curve,akcp-8-7-diagonal-hilbert}
For a flat finitely presented locally projective family of integral curves of
common genus $p>0$, the first Abel map
\[
\alpha_1:\Hilb^1_{X/S}\longrightarrow P_{-p}
\]
is a closed embedding, canonically identified with a closed embedding
$X\hookrightarrow P_{-p}$.  Its inverse image of the invertible Picard locus is
the smooth locus of $X/S$.
\end{theorem}
\begin{proof}
By properness and finite presentation of the Abel map, it suffices to show that
each geometric fiber is empty or one reduced point.  Over an algebraically
closed field, two points in one fiber give isomorphic maximal ideals
$\mathfrak m_x\simeq\mathfrak m_y$.  Since they are rank-one torsion-free,
the isomorphism is multiplication by a rational function $g$; distinctness
forces both ideals invertible and $1,g$ generate $\mathfrak m_x^{-1}$.
They therefore define $h:X\to\PP^1$.  The only pole of $g$ is the simple pole
at $x$, so $g$ generates the function field and $h$ is birational.  A proper
birational map from the integral curve to $\PP^1$ is an isomorphism, forcing
$p=0$, contradiction.  Thus fibers are single reduced points and the map is
a closed embedding.  The diagonal identification is \cref{akcp-8-7-diagonal-hilbert};
a maximal ideal is invertible exactly at a smooth point, proving the final claim.
\end{proof}
